\documentclass[a4paper]{article}
\setlength{\topmargin}{-1.0in}
\setlength{\oddsidemargin}{-0.2in}
\setlength{\evensidemargin}{0in}
\setlength{\textheight}{10.5in}
\setlength{\textwidth}{6.5in}
\usepackage{enumitem}
\usepackage{amsmath}
\usepackage{hyperref}
\usepackage{amssymb}
\usepackage[dvipsnames] {xcolor}
\usepackage{mathpartir}
\usepackage{algorithm}
\usepackage{algpseudocode}
\usepackage{csquotes}

\hbadness=10000

\hypersetup{
    colorlinks=true,
    linkcolor=blue,
    filecolor=magenta,      
    urlcolor=cyan,
    pdftitle={Assignment 1},
    pdfpagemode=FullScreen,
    }
\def\endproofmark{$\Box$}
\newenvironment{proof}{\par{\bf Proof}:}{\endproofmark\smallskip}
\begin{document}
\begin{center}
{\large \bf \color{red}  Department of Computer Science} \\
{\large \bf \color{red}  Ashoka University} \\

\vspace{0.1in}

{\large \bf \color{blue} Advanced Algorithms}

\vspace{0.05in}

    { \bf \color{YellowOrange} Assignment 1}
\end{center}
\medskip

{\textbf{Collaborators:} None} \hfill {\textbf{Name: Rushil Gupta} }

\bigskip
\hrule


% Begin your assignment here %
\section*{Problem 2}
Consider the bichromatic coloring problem: Given a set $S$ of $n$ elements, and subsets $S_1, S_2, \ldots, S_k$ of $S$, find a 2-coloring of the elements in $S$ such that every $S_i$ contains at least one red and at least one green element. \\

\noindent We will analyse the randomized algorithm that independently colors each element of $S$ red or green with probability $1/2$ under the Las Vegas, Monte Carlo and derandomized settings. \\

\subsection*{Las Vegas}
\noindent The Las Vegas algorithm will keep randomly coloring the elements of $S$ until a valid coloring is found. We are interested in the expected number of iterations before a valid coloring is found. Since each trial is independent, this problem reduced to finding the success probability of a single trial. \\

\noindent Let $A_i$ be the event that the subset $S_i$ is not bichromatic, i.e. a failiure event. Then, $Pr(A_i) = 2^{1-|S_i|}$, since there are two ways in which $S_i$ can be monochromatic (all red or all green) and $2^{|S_i|}$ total ways to color the elements of $S_i$. \\

\noindent Let $A$ be the event that the coloring is not valid, i.e. at least one of the subsets is not bichromatic. Then, by union bound, we have
\[Pr(A) \leq \sum_{i=1}^{k} Pr(A_i) = \sum_{i=1}^{k} 2^{1-|S_i|}\] \\

\noindent Therefore, the probability of success in a single trial is
\[Pr(\text{success}) = Pr(A^c) = 1 - Pr(A) \geq 1 - \sum_{i=1}^{k} 2^{1-|S_i|}\]

\noindent A remark here is that for this to be a useful lower bound on the failure probability, it suffices to have $\sum_{i=1}^{k} 2^{1-|S_i|} < 1$. \\

\noindent Assuming this is true, the expected number of trials before a success is
\[\mathbb{E}[\text{\# trials}] = \frac{1}{Pr(\text{success})} \leq \frac{1}{1 - \sum_{i=1}^{k} 2^{1-|S_i|}}\]

\noindent The cost of verifying if a coloring is valid is $O(N)$, where $N = \sum_{i=1}^{k} |S_i|$. Therefore, an upper bound on the expected running time of the Las Vegas algorithm is:

\[O\left(\frac{N}{1 - \sum_{i=1}^{k} 2^{1-|S_i|}}\right)\]


\vspace{0.2in}
\subsection*{Monte Carlo}
\noindent The same argument as above shows that the probability of failure in a single trial is at most $\sum_{i=1}^{k} 2^{1-|S_i|}$. Therefore, if we run the algorithm for $t$ trials, the probability of failure is at most:
\[Pr(\text{fail in all } t \text{ trials}) \leq \left(\sum_{i=1}^{k} 2^{1-|S_i|}\right)^t\]

\noindent The running time of the Monte Carlo algorithm is $O(Nt)$, where $N = \sum_{i=1}^{k} |S_i|$. Therefore, for any $\delta > 0$, if we want the probability of failure to be at most $\delta$, we can set
\[t \geq \frac{\log \delta}{\log \left(\sum_{i=1}^{k} 2^{1-|S_i|}\right)}\]

\noindent The running time of the algorithm is independent of the success probability, and is:
\[O\left(N \cdot t\right)\]

\noindent where $t$ can be chosen based on the desired probability of failure $\delta$.

\vspace{0.2in}
\subsection*{Derandomized}
Consider a graph-like data structure where each element in $S$ is a vertex, and each subset $S_i$ is also a vertex. There is an edge between an element vertex and a subset vertex if the element is in the subset. \\

\noindent This data structure maintains a partial coloring of the elements in $S$, and for each $S_i$, we know the number of colored elements of each color (keep a counter on the subset vertices). This also allows us to know all the subsets a given element belongs to. \\

\noindent We will iteratively color the elements of $S$ one by one, and at each step, we find the failure probability condition on coloring the current element red or green, and choose the color that leads to the lower failure probability. The proof of correctness for this given that the initial failure probability is less than 1 was done in the lecture.\\

\noindent By the proof covered in the lecture, knowing which case we lie in (1, 2 or 3). Knowing the number of uncolored elements in each subset, we can compute the failure probability of the current partial coloring conditioning on coloring the current element red or green. \\

\noindent It is important to see that this is a constant amount of work for each subset. So, at each step, the work done is proportional to the number of subsets the current element belongs to. This gives us the time complexity of the derandomized algorithm as:

\[O(N) = O\left(\sum_{i=1}^{n} \text{\# of times element } i \text{ appears in subsets}\right) = O\left(\sum_{i=1}^{k} |S_i|\right)\]\\


\noindent\textbf{Remark:} This discussion and analysis of the problem and 3 solutions did not need the subsets to be of the same size. The analysis holds for arbitrary subset sizes.


\newpage \section*{Problem 3}
Consider the $0$--$1$ Knapsack problem: we have a capacity $C$ and items with weight $w_i$ and profit $p_i$. We look at a simple greedy algorithm based on the profit density $p_i / w_i$. Sort all items in non-increasing order of $p_i / w_i$: $x_1, x_2, \ldots$. Take items in this order while they fit. Suppose we successfully take $x_1,\ldots,x_k$, but the next item $x_{k+1}$ does not fit. Let
\[
G_1 = \sum_{i=1}^k p_i \qquad \text{and} \qquad G_2 = p_{k+1}.
\]
We output
\[
G = \max\{G_1,\; G_2\}.
\]\\

\noindent \textbf{Claim:} This value $G$ is at least half of the optimal $0$--$1$ knapsack profit. \\

\noindent Consider the fractional version of knapsack where we are allowed to take a fraction of one item. The greedy-by-density method is optimal for the fractional problem: it takes all of $x_1,\ldots,x_k$ and then as much as possible of $x_{k+1}$. Let $\text{OPT}$ be the optimal 0-1 knapsack profit and let $\text{OPT}_{\text{frac}}$ be the fractional profit of this greedy solution. We have $\text{OPT} \le \text{OPT}_{\text{frac}}$ because allowing fractions cannot reduce the profit.\\

\noindent The fractional value is
\[
	\text{OPT}_{\text{frac}} = \sum_{i=1}^k p_i + \Big(C - \sum_{i=1}^k w_i\Big) \cdot \frac{p_{k+1}}{w_{k+1}}.
\]
Since $x_{k+1}$ did not fit completely, we have $C - \sum_{i=1}^k w_i < w_{k+1}$. Therefore
\[
	\text{OPT}_{\text{frac}} < \sum_{i=1}^k p_i + p_{k+1} = G_1 + G_2.
\]
Combining gives
\[
	\text{OPT} \le \text{OPT}_{\text{frac}} < G_1 + G_2.
\]\\
Thus at least one of $G_1$ or $G_2$ is at least half of $\text{OPT}$, i.e.
\[
\max\{G_1, G_2\} \ge \tfrac{1}{2}\, \text{OPT}.
\]\\

\noindent This is precisely saying that this modified greedy algorithm gives at least half of the optimal profit. \\

\newpage \section*{Problem 10}

We have a bin with $n$ balls, where $n$ is even. Exactly $n/2$ are red and $n/2$ are black. We draw balls one by one until we get a red ball. We compare two schemes: i) with replacement, and ii) without replacement.\\

\noindent Let $T$ be the draw number on which we first see a red ball.

\subsection*{With replacement}
Each draw is an independent Bernoulli trial with success probability $p = 1/2$ (success = red). Therefore
\[Pr(T = k) = (1-p)^{k-1} p = (1/2)^k, \qquad k = 1,2,3,\ldots\]
and
\[Pr(T > k) = (1-p)^k = (1/2)^k.\]
This is the standard geometric distribution. Its mean is
\[\mathbb{E}[T] = \frac{1}{p} = 2\]

\subsection*{Without replacement}
Here, the analysis is a bit more tedious, but it is intuitive to think why this should help: after seeing a black ball, there is are more red balls than black balls, so the chance of seeing a red ball in the next draw increases.\\

\noindent We will do the following analysis. Consider a sequence of $n$ balls, with $n/2$ red and $n/2$ black, with all such sequences are equally likely. Place a red ball at the start and end of the sequence. Now, the number of draws until we see a red ball is the position of the first red ball in this modified sequence. So, we want to know the expected length of the first segment of black balls in this modified sequence. \\

\noindent Now, we have $\frac{n}{2} + 2$ red balls (ends of contiguious segments of black balls) and $\frac{n}{2} + 1$ gaps between them (segments of black balls). Then, the expected length of these gaps are the same due to symmetry and the fact that we are considering all sequences equally likely. Since we have $n + 1$ balls in the sequence (we are not counting the red ball at the start since we think of that as the "0th" draw), the expected length of each gap is:

\[\mathbb{E}[T] = \frac{n + 1}{\frac{n}{2} + 1}\]

\noindent Clearly this is less than 2 for all $n \geq 0$, and tends to 2 as $n \to \infty$. Thus, the expected number of draws until we see a red ball is less when we draw without replacement, as intuition suggested. \\


\newpage
\section*{Problem 12}
We will analyze the space usage and search time of a randomized skip list with $n$ elements where each element is independently promoted to the next higher level with probability $1/2$.\\

\noindent Let $S(n)$ be the total space (number of node copies across all levels) and $T(n)$ be the search time for any fixed element.

\subsection*{Search Time}
We want to show that for some constant $c > 0$ and any $\alpha \geq 1$:
\[Pr(T(n) \geq \alpha \cdot c \cdot \log n) \leq \frac{1}{n^\alpha}\]

\noindent The search process moves right within levels and down between levels. Let $H$ be the maximum height of the skip list. The total search cost is proportional to $H$ plus the number of horizontal moves, so we will bound both these quantities separately. \\

\noindent The probability that any element reaches height $h$ is $2^{-h+1}$, since it is a geometric random variable. By union bound:
\[Pr(H \geq h) \leq n \cdot 2^{-h+1}\]

\noindent Setting $h = 2\alpha \log n$ gives:
\[Pr(H \geq 2\alpha \log n) \leq n \cdot 2^{-2\alpha \log n + 1} = \frac{2}{n^{2\alpha}} \leq \frac{1}{n^\alpha}\]

\noindent Now, note that at each level, the expected number of horizontal moves is constant (at most 2) due to the geometric distribution of gaps between promoted elements. So:
\[Pr(T(n) \geq c \alpha \log n) \leq \frac{1}{n^\alpha}\]

\subsection*{Space Bound}
We want to show that for any $a > 0$:
\[Pr(S(n) \geq 2n + an) \leq 2^{-a^2 n}\]

\noindent First, note that the expected space is $\mathbb{E}[S(n)] = 2n$ since each of the $n$ elements has expected height 2. Now, we know for $\delta \in [0,1]$:
\[Pr(S(n) \geq (1+\delta) \cdot 2n) \leq e^{-\frac{\delta^2}{3} \cdot 2n}\]

\noindent Setting $\delta = a/2$ gives:
\[Pr(S(n) \geq 2n + an) \leq e^{-\frac{a^2}{12} \cdot 2n} = e^{-\frac{a^2}{6} n} \leq 2^{-a^2 n}\]
\end{document}